%% kapa: exam=1 round=24 year=2013 period=1 form=A subject=civil_law subject_ko=민법 seq=20 number=20
%% source: Q-Net 감정평가사 기출문제 (cst00309) artlSeq=5073802 file=extracted/1008/5073802/2013년 제24회 감정평가사 1차시험 기출문제/1교시(A형).hwp
%% conversion: kordoc 4.24.51 -> tools/gen1.py (auto); review=proofread-sonnet; flags=none
\begin{question}{20}
물권적 청구권에 관한 설명으로 옳은 것은? (다툼이 있으면 판례에 의함)
\begin{choices}
  \item 점유물반환청구권은 점유를 침탈당한 날로부터 1년 내에 행사하여야 한다.
  \item 양자간 등기명의신탁에서 명의수탁자가 신탁부동산을 처분하여 제3취득자가 유효하게 소유권을 취득하였는데, 그 후 명의수탁자가 우연히 신탁부동산의 소유권을 다시 취득하였다면 명의신탁자의 소유권에 기한 물권적 청구권이 인정된다.
  \item 제3자가 저당권설정자인 소유자의 의사에 반하여 저당권이 설정된 건물의 점유를 침탈한 경우 저당권자는 자신에게 건물의 반환을 청구할 수 있다.
  \item 토지소유자는 그의 토지 위에 무단으로 건축된 건물을 임차하여 점유하고 있는 자를 상대로 건물철거를 청구할 수 있다.
  \item 유효하게 부동산을 명의신탁한 자는 자신이 직접 제3자에게 물권적 청구권을 행사하여 신탁재산에 대한 침해배제를 구할 수 있다.
\end{choices}
\end{question}
\answer{1}
